\section{Discussion and conclusion}\label{sec:discussion}

The experiments support a narrower and more useful conclusion than ``RL works'' or ``RL does
not work.'' RL is warranted only when a simulation exposes observable, state-contingent,
resource-adjusted value. Physical authority alone is insufficient: an actuator may change a
trajectory while the same posture remains optimal. Perfect-information headroom is also
insufficient: an oracle may exploit future events unavailable to an operator. Finally, action
prediction is insufficient: classification accuracy does not weight mistakes by their
sequential regret.

Comparator design is part of the estimand, not a robustness appendix. The Track B reversal
shows that a stable held-out gain against a restricted frontier can disappear when constants
receive the learner's complete rights. This does not make the earlier computation erroneous;
it changes what it identifies. The appropriate claim is that downstream and upstream rights
matter and restricted comparator families can exaggerate adaptive value. Studies should state
which rights are shared, which are fixed, and whether resources are equal, dominated, or merely
priced in a reward.

The null results are cumulative evidence rather than failed attempts. Allocation failed because
automatic reallocation and symmetry removed durable opportunity cost. Finite convoy control
created clairvoyant headroom but did not yield observable conversion. A readiness portfolio
created specialized actions under a fixed budget but still failed rollout conversion. Spatial
tempo and advance signals did not beat blind alternation after ledger correction. Together they
map a causal chain from physical sensitivity through information and policy realization.

Metric triangulation adds a second warning. Order continuity, ration mass, spatial fairness,
and multidimensional aggregate indices answer different managerial questions. Reporting them
together exposes trade-offs; selecting the most favorable one after evaluation converts
construct choice into model selection. Cobb--Douglas remains valuable for decomposing
inventory, backlog, time, capacity, and resource preferences, but its weights and variables
must be validated for the decision context.

The principal limitation is scope. The findings arise from one thesis-grounded reconstruction
and disclosed extensions. Program G is stylized rather than a full event-driven spatial port.
The experiments do not establish that all military logistics decisions lack adaptive value,
nor that a mathematically optimal partially observable policy cannot outperform the policies
tested. They establish where the tested contracts failed and which stronger claims are not
licensed. A final, separately preregistered information-sufficiency audit may distinguish
policy-class limitation from information limitation without changing Program G's verdict.

For practice, the framework provides a stopping discipline. Before funding learner training,
an organization can branch actions under exact replay, enumerate strong static doctrines,
measure action-ranking diversity, fit a transparent observable policy, and audit resource use.
If these inexpensive stages fail, a neural policy is unlikely to create value. If they pass,
RL enters as one comparator among trees, heuristics, belief-aware planning, and model-predictive
control rather than as the default solution.

In conclusion, the contribution is an eligibility framework and an auditable empirical
demonstration of its necessity. Across the MFSC decision families, apparent adaptive gains are
absorbed by same-contract statics, disappear at observable conversion, or depend on the chosen
resilience functional. The result is not a rejection of reinforcement learning. It is a rule
for when training an agent constitutes a scientifically identified experiment.
